\scanpage{30}
% sheet: reused_typed
% hand: typed
% L1 | top centre, in reading orientation
\mspage{-2.17-}

% L2–5 | top, end of the preceding proof
\[
\begin{aligned}
\lambda(\alpha)\cdot\gamma\cdot\lambda(\alpha^{-1})
 &=\handwritten{\chi}\bigl(\lambda(\alpha^{-1})\bigr)\cdot\gamma\\
 &=\alpha^{-\handwritten{\bar{\chi}}(\handwritten{\bar{\lambda}})}\cdot\gamma
\end{aligned}
\]
for $\gamma\in(\handwritten{\mathbb G}_a)_k$. Therefore $\lambda(\alpha)\cdot\gamma\cdot\lambda(\alpha^{-1})$ has a finite specialization when $\alpha\handwritten{\longrightarrow}0$, i.e. $\handwritten{\mathbb G}_a\handwritten{\subset}P(\lambda)$. Therefore $B\handwritten{\subset}P(\lambda)$. \uline{QED}

% L6–9 | upper middle
\uline{Proposition} \uline{2.7}: Let a reductive group $G$ act via $\sigma$ on a scheme $X$, proper over $k$. Then for all $x\in X_k$, $L\in\mathrm{Pic}^{G}(X)$, and 1-PS's $\lambda$ of $G$,
\[
\mu^{L}(x,\lambda)=\mu^{L}(x,\gamma^{-1}\cdot\lambda\cdot\gamma)
\]
\margin{\handwritten{\unsure{9}}}
if $\gamma\in P(\lambda)_k$.

% L10–14 | middle, start of proof and calculation
\uline{Proof}: Let $y$ be the specialization of $\sigma(\lambda(\alpha),x)$ as $\alpha\handwritten{\longrightarrow}0$. Suppose we write $y=\lim_{\handwritten{\alpha\to 0}}\sigma(\lambda(\alpha),x)$. In this notation, we calculate:
\[
\begin{aligned}
&\lim_{\handwritten{\alpha\to 0}}\sigma(\gamma^{-1}\cdot\lambda(\alpha)\cdot\gamma,x)\\
&\qquad=\lim_{\handwritten{\alpha\to 0}}\sigma\handwritten{\bigl(}\gamma^{-1}\cdot\lambda(\alpha)\cdot\gamma\cdot\lambda(\alpha^{-1}),\ \sigma(\lambda(\alpha),x)\handwritten{\bigr)}\\
&\qquad=\sigma\handwritten{\bigl(}\gamma^{-1}\cdot\lim_{\handwritten{\alpha\to 0}}\lambda(\alpha)\cdot\gamma\cdot\lambda(\alpha^{-1}),\ y\handwritten{\bigr)}
\end{aligned}
\]

% L15–19 | lower
Let $\lim_{\handwritten{\alpha\to 0}}\lambda(\alpha)\cdot\gamma\cdot\lambda(\alpha^{-1})=\gamma'$. According to Proposition 2.6, $\gamma'$ centralizes $\lambda$. Now $\mu(x,\gamma^{-1}\cdot\lambda\cdot\gamma)$ is calculated by looking at $\lim_{\handwritten{\alpha\to 0}}\sigma(\gamma^{-1}\cdot\lambda(\alpha)\cdot\gamma,x)$. Therefore:
\[
\mu(x,\gamma^{-1}\cdot\lambda\cdot\gamma)=\mu(\sigma(\gamma^{-1}\cdot\gamma',y),\gamma^{-1}\cdot\lambda\cdot\gamma)
\]

By property (i) of $\mu$ (\S1 above), we have:
% archive sheet number: 29

